\begin{lemma}
\label{lemma-NAK-localization}
Let $R$ be a ring, let $S \subset R$ be a multiplicative subset,
let $I \subset R$ be an ideal, and let $M$ be a finite $R$-module.
If $x_1, \ldots, x_r \in M$ generate $S^{-1}(M/IM)$
as an $S^{-1}(R/I)$-module, then there exists an $f \in S + I$
such that $x_1, \ldots, x_r$ generate $M_f$ as an
$R_f$-module.\footnote{Special cases: (I) $I = 0$. The lemma says
if $x_1, \ldots, x_r$ generate $S^{-1}M$, then $x_1, \ldots, x_r$
generate $M_f$ for some $f \in S$. (II) $I = \mathfrak p$ is
a prime ideal and $S = R \setminus \mathfrak p$. The lemma says if
$x_1, \ldots, x_r$ generate $M \otimes_R \kappa(\mathfrak p)$
then $x_1, \ldots, x_r$ generate $M_f$ for some
$f \in R$, $f \not \in \mathfrak p$.}
\end{lemma}

\begin{proof}
Special case $I = 0$. Let $y_1, \ldots, y_s$ be generators for $M$ over $R$.
Since $S^{-1}M$ is generated by $x_1, \ldots, x_r$, for each $i$
we can write $y_i = \sum (a_{ij}/s_{ij})x_j$ in $S^{-1}M$
for some $a_{ij} \in R$ and $s_{ij} \in S$. Multiplying by the product
$s=\prod_{k=1}^{s_0}\prod_{\ell=1}^r s_{k\ell}\in S$,
where $s_0$ denotes the number of chosen generators $y_i$, we see that
$sy_i = \sum_{j=1}^r a'_{ij}x_j$ in $S^{-1}M$, with
$$
a'_{ij}=a_{ij}
\prod_{\substack{1\leq k\leq s_0,\ 1\leq\ell\leq r\\(k,\ell)\ne(i,j)}}
s_{k\ell}\in R.
$$
Empty products are $1$ and empty sums are $0$.
This in turn means there exist $t_i \in S$ such that
$t_isy_i = \sum t_ia'_{ij}x_j$ in $M$. Thus if $t \in S$
is the product of the $t_i$, put
$t^{(i)}=\prod_{k\ne i}t_k$, so $t=t_it^{(i)}$.
Multiplying the equality in $M$ by $t^{(i)}$ gives
$sty_i=\sum_{j=1}^r t a'_{ij}x_j$.
In $M_{st}$ it follows that
$y_i=\sum_{j=1}^r(t a'_{ij}/(st))x_j$.
Thus $y_i$ is in the $R_{st}$-submodule generated by
$x_1,\ldots,x_r$. Hence $x_1, \ldots, x_r$ generate $M_{st}$.

\medskip\noindent
General case. By the special case, we can find an $s \in S$
such that $x_1, \ldots, x_r$ generate $(M/IM)_s$ over $(R/I)_s$.
By Lemma \ref{lemma-NAK} we can find a $g \in 1 + I_s \subset R_s$
such that $x_1, \ldots, x_r$ generate $(M_s)_g$ over $(R_s)_g$.
Write $g=1+i/s'$ with $i\in I$ and $s'=s^n$ for some
$n\geq0$. Put $f=ss'+is=s(s'+i)$. Since $ss'=s^{n+1}\in S$
and $is\in I$, we have $f\in S+I$.
In $(R_s)_g$, the image of $f$ equals $ss'g$ and is a unit.
Conversely, in $R_f$ the element $s$ has inverse $(s'+i)/f$,
and the image of $g$ has inverse $ss'/f$:
$$
s\frac{s'+i}{f}=1,\qquad
\left(1+\frac{i}{s'}\right)\frac{ss'}{f}=1.
$$
The two universal properties therefore give mutually inverse maps
$R_f\longleftrightarrow(R_s)_g$ which agree with the original
maps from $R$. More explicitly, the first sends $r/f^k$ to
$(r/1)/(ss'g)^k$. The reverse map sends
$r/s^a$ to $r(s'+i)^a/f^a$ and $g^{-1}$ to $ss'/f$.
Their composites fix the image of every $r\in R$ and each inverse
adjoined in the respective localizations, so both composites are identities.
The corresponding module maps are also inverse: they send
$m/f^k$ to $(m/1)/(ss'g)^k$ and send
$(m/s^a)/g^b$ to $m(s'+i)^a(ss')^b/f^{a+b}$.
They preserve the images of all the chosen $x_j$.
Consequently those images generate $M_f$ over $R_f$, as required.
These localization identities remain valid when a localization is
the zero ring and when the original lists of generators are empty.
\end{proof}